\documentclass[a4paper,10pt]{article}
\usepackage[francais]{babel}
\usepackage{url}

\selectlanguage{francais}
\input{letter-icube-uds.sty}

\selectlanguage{francais}
\begin{document}

\name{St\'ephane \textsc{Genaud}}
\fonction{Professeur des Universités}
\telephoneillkirch{+33 (0)3688 54542}
\email{genaud@unistra.fr}


\begin{letter}[Recommandation d'Etienne Michon pour le poste de Distributed application developer]{à}{Responsables recrutement pour France Life Imaging}

  Je  connais  Etienne  depuis  qu'il  a intégré  en  septembe  2010  le  Master
  informatique RISE de  l'Université de Strasbourg. Etienne,  après avoir achevé
  un cycle ingénieur, a choisi ce  Master orienté réseaux et systèmes distribués
  et parallèles  car cette formation lui  donnait la possibilité de  réaliser un
  stage de recherche dans le domaine. Etienne a effectivement choisi un stage de
  recherche de Master, co-encadré par Julien Gossa  et moi même, sur le thème de
  l'allocation et l'ordonnancement de ressources dans les clouds IaaS.

  À l'issue  de son Master recherche,  Etienne a obtenu un  financement de thèse
  lui permettant poursuivre son travail sur  le même thème. Etienne a pendant sa
  thèse beaucoup contribué a mettre au point notre logiciel \emph{Schlouder} qui
  sert   à  l'équipe   d'outil   scientifique  pour   tester   et  valider   nos
  propositions. Ce logiciel  est un courtier côté-client de  ressources de cloud
  IaaS.  Il  accepte  en  entrée des  stratégies  globales  d'ordonnancement  et
  d'allocation dynamiques multi-critères de  tâches indépendantes ou de workflow
  (souvent rapidité d'exécution versus le  coût économique). Cet outil permet de
  mener des  expériences dans  des environnements  réels de  clouds (s'interface
  avec OpenStack ou EC2).

  Sur  la dernière  partie de  sa thèse,  Etienne travaille  également au  volet
  suivant de  nos recherches  qui consistent  à proposer  une prédiction  des du
  résultats  des  stratégies  d'ordonnancement  choisies par  le  client.  Cette
  prédiction est  obtenue à partir  d'un simulation à évènements  discrets basée
  sur le logiciel \textsc{SimGrid}.

  Etienne soutiendra sa thèse de doctorat le  5 juin 2015. Son travail durant sa
  thèse a fait avancer de manière considérable notre réflexion, en particulier à
  travers les expériences scientifique menées par Etienne. Pour mettre en oeuvre
 les idées générales de l'équipe, Etienne a fourni un travail considérable et
 rigoureux de développement et d'analyse qui ont amené les réalisations que nous
décrivons dans un article de revue actuellement en soumission. Il a été amené à
transposer des applications scientifiques (en particulier une d'imagerie
astronomique, une en protéomique) de façon à les exécuter dans notre
environnement distribué. 

D'un point de  vue humain, Etienne a sans conteste  de grandes qualités d'écoute
et de rigueur  qui ont été extrêment bénéfiques dans  notre travail d'équipe. Il
sait prendre  en compte  les demandes  et les interrogations  et les  traiter de
manière systématique. De plus, Etienne a  un contact humain très agréable, d'une
humeur toujours égale même dans la difficuluté, qui rendent faciles les relations
de travail.
Enfin, je souligne les très grandes qualité techniques d'Etienne en matière
d'ingénierie informatique. Il a démontré lors de sa thèse qu'il était capable de
s'adapter très rapidement  à de nombreuses technologies ou langages auxquels
nous avons été confrontés dans nos développements.

Pour toutes ces raisons, je recommande Etienne Michon au plus haut point pour le
poste  de développeur  d'application distribuées,  impliquant de  la gestion  de
données  distribuées. Je  crois  qu'Etienne  à une  expérience  et des  qualités
très intéressantes vis à vis des exigences de ce travail.

\end{letter}
\begin{flushright}
\includegraphics[width=.20\textwidth]{signgenaud.jpg}
\end{flushright}
\end{document}
